\documentclass[11pt]{article}
\usepackage[margin=1in]{geometry}
\usepackage{amsmath,amssymb,amsthm}

\newtheorem{proposition}{Proposition}

\title{Early-Growth Thresholds for the SIS Attendance Model}
\author{Le Minh Ngo}
\date{}

\begin{document}
\maketitle

\paragraph{Setup.}
Let $A$ be the symmetric adjacency matrix of an undirected contact graph.
For attendance vector $d_t\in\{0,1\}^n$, set
$D_t=\operatorname{diag}(d_t)$ and $B_t=D_tAD_t$. Let $X_i(t)$ indicate
infection at the start of day $t$, let $p_i(t)=\Pr(X_i(t)=1)$, and put
$a=1-\gamma$, where $0<\gamma<1$. We use the rare-infection approximation
in which probabilities of simultaneous infections of two specified
individuals are second-order terms.

\begin{proposition}[Linearized 14-day growth]
To first order in the infection probabilities,
\begin{equation}
    p_{t+1}=M_t(\beta)p_t+O(\|p_t\|^2),
    \qquad M_t(\beta)=aI+\beta B_t.
    \label{eq:linear}
\end{equation}
For a 14-day attendance cycle, define
\begin{equation}
    P_{14}(\beta)=M_{13}(\beta)M_{12}(\beta)\cdots M_0(\beta),
    \qquad g_{14}(\beta)=\rho(P_{14}(\beta)).
    \label{eq:cycle}
\end{equation}
The linearized threshold $\beta_c$ satisfies $g_{14}(\beta_c)=1$.
\end{proposition}

\begin{proof}
Conditional on the start-of-day state, the SIS rule gives
\begin{equation}
\mathbb E[X_i(t+1)\mid X(t)]
=aX_i(t)+(1-X_i(t))
\left[1-(1-\beta)^{\sum_j(B_t)_{ij}X_j(t)}\right].
\end{equation}
With zero infected neighbors the bracket is zero, and with one it equals
$\beta$. Events with at least two infected neighbors and products
$X_i(t)X_j(t)$ contribute second-order terms under the stated
approximation. Taking expectations therefore yields
\[
p_i(t+1)=ap_i(t)+\beta\sum_j(B_t)_{ij}p_j(t)
+O(\|p_t\|^2),
\]
which is equation~\eqref{eq:linear}. Repeated substitution gives
$p_1\simeq M_0p_0$, $p_2\simeq M_1M_0p_0$, and
$p_{14}\simeq P_{14}p_0$ with the order in equation~\eqref{eq:cycle}.
The dominant per-cycle multiplier is the spectral radius $\rho(P_{14})$.
\end{proof}

\begin{proposition}[Schedule-specific thresholds]
Let $\lambda_F=\lambda_{\max}(A)$. Let $A_{00}$ and $A_{11}$ be the
adjacency matrices induced by the two rotating groups, and define
\[
\lambda_R=\max\{\lambda_{\max}(A_{00}),
                 \lambda_{\max}(A_{11})\}.
\]
Let $A_H$ be the graph obtained by permanently retaining a random half
of the original edges in the On-Off control, and put
$\lambda_H=\lambda_{\max}(A_H)$. For positive relevant eigenvalues,
\begin{align}
\beta_c^{\mathrm{Full}}&=\frac{\gamma}{\lambda_F},
    \label{eq:full}\\
\beta_c^{\mathrm{OnOff}}&=\frac{a^{-1}-a}{\lambda_F},
    \label{eq:onoff}\\
\beta_c^{\mathrm{Rotating}}&=\frac{a^{-1}-a}{\lambda_R},
    \label{eq:rotating}\\
\beta_c^{\mathrm{HalfEdge}}&=\frac{a^{-1}-a}{\lambda_H}.
    \label{eq:half}
\end{align}
The daily and weekly versions of each 50\%-attendance schedule have the
same linearized 14-day threshold. If a relevant eigenvalue is zero, its
corresponding threshold is $+\infty$.
\end{proposition}

\begin{proof}
Because $A$ is symmetric and entrywise nonnegative, the leading
eigenvalue of $aI+\beta A$ is $a+\beta\lambda_F$ for $\beta\geq0$.
Under Full attendance,
\[
g_{14}^{\mathrm{Full}}(\beta)
=\rho\bigl((aI+\beta A)^{14}\bigr)
=(a+\beta\lambda_F)^{14}.
\]
Equating this expression to one gives
$a+\beta_c\lambda_F=1$, hence
$\beta_c=(1-a)/\lambda_F=\gamma/\lambda_F$.

For On-Off, any complete 14-day cycle has seven on days with matrix
$aI+\beta A$ and seven off days with matrix $aI$. The off-day matrix
commutes with the on-day matrix, so either switching frequency has
\[
P_{14}^{\mathrm{OnOff}}(\beta)
=(aI)^7(aI+\beta A)^7,
\qquad
g_{14}^{\mathrm{OnOff}}(\beta)
=\bigl[a(a+\beta\lambda_F)\bigr]^7.
\]
Thus $a(a+\beta_c\lambda_F)=1$, or
$\beta_c=(a^{-1}-a)/\lambda_F$.

For Rotating, reorder vertices by group. The two daily matrices are
\[
M^{(0)}
=\begin{pmatrix}
aI_0+\beta A_{00}&0\\0&aI_1
\end{pmatrix},
\qquad
M^{(1)}
=\begin{pmatrix}
aI_0&0\\0&aI_1+\beta A_{11}
\end{pmatrix}.
\]
They commute because they act on disjoint blocks. Each group attends
seven days and stays home seven days, giving
\[
P_{14}^{\mathrm{Rotating}}(\beta)
=\begin{pmatrix}
a^7(aI_0+\beta A_{00})^7&0\\
0&a^7(aI_1+\beta A_{11})^7
\end{pmatrix}.
\]
Therefore
$g_{14}^{\mathrm{Rotating}}(\beta)
=\bigl[a(a+\beta\lambda_R)\bigr]^7$,
and equation~\eqref{eq:rotating} follows. Finally, the half-edge
control has the On-Off attendance matrices with $A$ replaced by $A_H$,
which yields equation~\eqref{eq:half}.
\end{proof}

\paragraph{Two-vertex check.}
For one edge connecting two vertices, $\lambda_F=1$. With
$\gamma=0.2$, $a=0.8$ and $a^{-1}-a=1.25-0.8=0.45$.
Equation~\eqref{eq:onoff} gives $\beta_c^{\mathrm{OnOff}}=0.45$.
For the 200-person networks, the same formula divides $0.45$ by each
world's much larger $\lambda_F$, giving a mean threshold near $0.05$.

\end{document}
